\section{Comparison with base papers and state of the art}\label{sec:base}
The methodological antecedents include domain-independent alert aggregation, multi-stage causal reconstruction, provenance-based investigation, and incident triage \cite{landauer2022,grain2025,holmes2019,sadlek2025,crosshound2026}. Table~\ref{tab:base_papers} records their tasks and published evaluation settings. The present method addresses a narrower decision: whether a particular cross-fragment union is justified by three original, timely, independently rooted observations, followed by ranking of the resulting partition. Its 80\% complete-and-pure incident Recall@20 is not the same metric as attack-step coverage, alert reduction, sequence classification, or severity discrimination.

\begin{table}[t]\centering\footnotesize
\caption{Base-paper comparison. Published outcomes are marked ``reported'' and are never treated as results on this paper's cohort. A numeric cross-paper improvement is undefined without a matched evaluation protocol.}\label{tab:base_papers}
\begin{tabularx}{\textwidth}{@{}p{.16\textwidth}p{.21\textwidth}p{.21\textwidth}Y@{}}\toprule
Study & Problem and method & Evaluation setting and reported result & Difference and comparability\\\midrule
Landauer et al. \cite{landauer2022} & Domain-independent learned alert aggregation & Semi-structured alert datasets; aggregation quality reported in source & Different alert definitions and grouping objective; no same-cohort reproduction\\
GRAIN \cite{grain2025} & Causal multi-step reconstruction with graph learning and reinforcement learning & Public attack-scenario datasets; graph/path reconstruction reported in source & Different ground truth and output; no same-cohort reproduction\\
HOLMES \cite{holmes2019} & Suspicious information-flow correlation for APT detection & System-level provenance evaluation reported in source & Host audit-event flow differs from ranked alert partitions\\
Sadlek et al. \cite{sadlek2025} & Kill-chain graph severity triage & Testbed sequence classification; reported ROC AUC 0.95 & AUC for malicious sequences is not Incident Recall@20\\
CrossHound \cite{crosshound2026} & Host/network provenance correlation and attack reconstruction & OpTC/LANL; reported 97\% attack-step coverage on OpTC & Attack-step coverage is not 80\% complete-and-pure top-20 incident retrieval\\
This study & Evidence-gated incident union plus deterministic rank; bounded C/P/V role framework & 15 synthetic queues; 260/300 R@20 and zero observed false pairs; LLM efficacy unmeasured & Fresh paired baseline only within synthetic cohort\\\bottomrule
\end{tabularx}\end{table}

The only valid numerical improvement in this manuscript is on the paired synthetic holdout: the candidate increases Recall@20 by $10/300=3.33$ percentage points (relative change $10/250=4.0\%$), recoverability by $18/300=6.00$ percentage points (relative $18/263\approx6.84\%$), and P@10 by $6/150=4.00$ percentage points (relative $6/107\approx5.61\%$). Both arms use identical eligible alerts and the same evaluator. No base paper was reimplemented on this cohort; their published numbers cannot be placed in a shared performance ranking. Their methodological differences motivate separate provenance, arrival-time and safety gates, while their public data do not supply the frozen verifier's complete input contract.
